\setcounter{aufgabe}{7}
\einheitenkopf{Einheit 1 von 5 · Häufigkeiten}

\begin{aufgabe}{Häufigkeiten -- Was ist das Ganze? Unterstreiche im Text die Zahl, die das Ganze angibt, und schreibe sie ins Feld. Rechne noch nicht.}
\begin{teile}
\teil In einer Woche fanden 18 Spiele statt. 7 davon gewann die Mannschaft. \leerfeld
\teil Von den 250 Besuchern kamen 90 mit dem Fahrrad. \leerfeld
\teil In einer Klasse sind 14 Mädchen und 12 Jungen. 9 von ihnen spielen ein Instrument. \leerfeld
\teil An der Umfrage nahmen 40 Personen teil. 6 nannten Tennis, 14 Fußball, die übrigen Handball. \leerfeld
\teil In einer Kiste liegen 12 rote, 8 blaue und 5 grüne Kugeln. Welchen Anteil haben die blauen? \leerfeld
\end{teile}
\end{aufgabe}

\begin{aufgabe}{Häufigkeiten -- auszählen. Zähle die Striche aus und schreibe die Anzahl in die rechte Spalte.}
\beispiel{\text{\strichliste{|||||\ ||}} \quad &\rightarrow\quad 7}

\sachtabelle{lcc}{Farbe & Strichliste & Anzahl}{%
a) Rot & \strichliste{|||||\ |||} & \leerzelle \\
b) Blau & \strichliste{|||||\ |||||\ ||} & \leerzelle \\
c) Grün & \strichliste{||||} & \leerzelle \\
d) Gelb & \strichliste{|||||\ |||||\ |||||\ |} & \leerzelle}

\begin{teile}
\setcounter{teil}{4}
\teil Beim Würfeln wurden diese Augenzahlen notiert (Urliste):
\par\smallskip
\noindent 3\; 5\; 1\; 3\; 6\; 3\; 2\; 5\; 3\; 1\; 6\; 4\; 3\; 2\; 5\; 3\; 6\; 1\; 4\; 3
\par\smallskip
Trage die Anzahlen in die Häufigkeitstabelle ein.
\end{teile}

\sachtabelle{lcccccc}{Augenzahl & 1 & 2 & 3 & 4 & 5 & 6}{%
Anzahl & \leerzelle & \leerzelle & \leerzelle & \leerzelle & \leerzelle & \leerzelle}
\end{aufgabe}

\begin{aufgabe}{Häufigkeiten -- rechnen. Gib bei a) bis e) die relative Häufigkeit als Bruch an.}
\beispiel{\text{9 von 12:}\quad h &= \tfrac{9}{12}}
\begin{teilezwei}
\tz 3 von 10 \leerfeld & \tz 7 von 20 \leerfeld \\
\tz 9 von 25 \leerfeld & \tz 4 von 50 \leerfeld \\
\end{teilezwei}
\begin{teile}
\teil 6 von 24 -- kürze den Bruch so weit wie möglich. \leerfeld
\end{teile}

Jetzt als Dezimalzahl oder in Prozent. Du darfst den Taschenrechner benutzen.
\begin{teile}
\teil 13 von 20 -- als Dezimalzahl. \feld{h}
\teil 17 von 40 -- als Dezimalzahl und in Prozent. \feld{h} \leerfeld[\%]
\teil 5 von 9 -- in Prozent, gerundet auf eine Stelle nach dem Komma. \leerfeld[\%]
\teil In einer Umfrage antworten $45\,\%$ „ja“ und $35\,\%$ „vielleicht“, die übrigen „nein“. Prüfe, ob die drei Anteile zusammen $100\,\%$ ergeben, und gib den fehlenden Anteil an. \leerfeld[\%]
\teil Lies im Diagramm ab, wie viele Jugendliche Tanzen nennen und wie viele insgesamt befragt wurden. Gib die relative Häufigkeit für Tanzen in Prozent an. \leerfeld[\%]
\end{teile}

\noindent\saeulen{Fußball/12,Tanzen/9,Rad/15,Judo/4}{18}{3}{Anzahl}

\begin{teile}
\teil Ein Verein hat 320 Mitglieder. $15\,\%$ von ihnen sind jünger als 12 Jahre. Wie viele Mitglieder sind das? \leerfeld
\teil Eine Mannschaft bestritt in einer Saison 40 Spiele. Sie gewann 26 Spiele, die übrigen verlor sie. Gib die relative Häufigkeit der Siege als Bruch und in Prozent an. (P10 2015)
\par \feld{h} \leerfeld[\%]
\end{teile}
\end{aufgabe}

\begin{aufgabe}{Häufigkeiten -- Fehler finden.}
In einer Klasse mit 32 Kindern haben 8 einen Hund. Mia gibt als relative Häufigkeit an:
\rechnung{h &= 8}
\begin{teile}
\teil Finde den Fehler, benenne ihn und schreibe die richtige Rechnung auf. \\[12pt]
\teil In einer Gruppe von 45 Personen tragen 18 eine Brille. Gib die relative Häufigkeit als gekürzten Bruch und in Prozent an. \par \feld{h} \leerfeld[\%]
\end{teile}
\end{aufgabe}

\begin{aufgabe}{Häufigkeiten -- Begründen. Rechne bei b) nicht mit dem Taschenrechner.}
\begin{teile}
\teil Erkläre, warum man zwei verschieden große Gruppen nur mit relativen Häufigkeiten vergleichen kann und nicht mit den Anzahlen. \\[12pt]
\teil In der 7a haben 12 von 24 Kindern ein Haustier, in der 7b sind es 15 von 30. In welcher Klasse ist der Anteil größer? Begründe. \\[12pt]
\end{teile}
\end{aufgabe}

\begin{aufgabe}{Häufigkeiten -- Anwendung. An zwei Schulen wurde gezählt, wie viele Kinder mit dem Rad kommen: an der Nordschule 84 von 240 Kindern, an der Südschule 78 von 195 Kindern.}
\begin{teile}
\teil Gib für beide Schulen die relative Häufigkeit in Prozent an.
\par Nordschule \leerfeld[\%] \quad Südschule \leerfeld[\%]
\teil An welcher Schule kommen mehr Kinder mit dem Rad? \leerfeld
\teil An welcher Schule ist der Anteil der Radfahrenden größer? \leerfeld
\teil Die Stadt baut zuerst an der Schule mit dem größeren Anteil neue Fahrradständer. Welche Schule ist das, und warum reicht die Anzahl aus b) für diese Entscheidung nicht aus? \\[12pt]
\end{teile}
\end{aufgabe}
